\section{Experiments}\label{sec:experiments}

We organize our experiments around the contributions from \cref{sec:introduction}: our GPP project results (\cref{sec:exp:gpp}), \texttt{Continual Harness} closing the gap to a hand-engineered harness (\cref{sec:exp:closegap}), showing improvements with reset-free experience that can bootstrap runs if one does choose to reset, and continual model-harness co-learning for open-source students (\cref{sec:exp:colearn}). 
\cref{sec:exp:mechanisms} attributes these gains to in-loop refinement on each of the harness components, and additional details are in the appendix.

\subsection{Setup}\label{sec:exp:setup}

\paragraph{Environments and metric.} We evaluate on Pok\'emon Red and Emerald, two RPGs in the same genre that differ in map layout, mechanics, and difficulty. We use the standardized milestone evaluation from the PokeAgent Challenge~\citep{karten2026pokeagent}. The primary metric is \textbf{cumulative button presses to milestone}.

\paragraph{Harness conditions.} $\mathcal{H}_{\min}$: frames, local text map, buttons, generic system prompt; no sub-agents, memory, or skills. $\mathcal{H}_\mathrm{expert}$: the hand-designed harness of PokeAgent~\citep{karten2026pokeagent} and fixed GPP harness with built sub-agents, A$^*$ pathfinding, type chart, damage calculator, and curated objectives. $\mathcal{H}_\mathrm{CH}$: starts from $\mathcal{H}_{\min}$ and refines during gameplay via \cref{fig:methodology}; three variants: \emph{from scratch}, \emph{bootstrap frozen} (loads a successful from-scratch run, refinement disabled), \emph{bootstrap updating} (same bootstrap, refinement continues).

\paragraph{Models and seeds.} we use Gemini 3 variants (Pro, Flash, Flash-Lite) across all harness conditions, and for open-source transfer (\cref{sec:exp:colearn}) we use Gemma-4 (E2B, E4B, 26B MoE, 31B dense). We use at least three seeds across all experiments. We report seed medians with per-seed traces at reduced opacity.

\subsection{Gemini Plays Pok\'emon completes multiple RPGs}\label{sec:exp:gpp}

Our GPP project ran Gemini models live through Pok\'emon Blue (May 2025), Yellow Legacy on hard mode (August 2025), and Crystal without a lost end-game battle (November 2025), making GPP the first AI system to complete multiple Pok\'emon RPGs. Since GPP used a mix of human designed and agent iterated harness, we highlight specific cases where we explored harness refinement over thousands of hours of gameplay.

\paragraph{Emergent Continual Harness behavior through skills.} Our Blue-era GPP harness relied on hand-authored specialists such as Pathfinder Agent and Boulder Puzzle Strategist. From Yellow Legacy onward we replaced these with general skills (\texttt{define\_agent}, \texttt{run\_code}, notepad edits) and let the model build its own harness. Unprompted behaviors included wrapping an \texttt{autopress\_buttons} sandbox loophole into a general \texttt{press\_sequence} primitive, developing named multi-stage battle strategies (``Operation Zombie Phoenix'' on Crystal's final Red fight), and authoring an explicit truth-table representation of the Goldenrod Underground switch puzzle in the notepad.

\paragraph{Quantitative harness growth.} \cref{fig:gpp_yellow} reports CRUD operations (creation, update, delete) on skill and sub-agent definitions across our Yellow Legacy run. Updates persist throughout the run rather than converging to a fixed harness, and concentrate on a small subset of navigation and battle components. \cref{fig:gpp_battle_complexity} reports structural metrics of one such component, the \texttt{battle\_strategist\_agent} prompt, across successive revisions during the Elite Four phase. The prompt cycles between growth and simplification, and undergoes a structural rewrite in which per-decision logic is absorbed into a \texttt{master\_battle\_agent} that dispatches to named sub-checks. Across both figures the process is the same: a small set of components is repeatedly updated and periodically rewritten. Quality is established separately by GPP's completion record. We generalize GPP's mixed methodology to create Continual Harness, which fully automates this process for all modules. Additional results in Appendix~\ref{app:gpp}.

\begin{figure}[t]
\centering
\includegraphics[width=\linewidth]{h0_gpp_yellow.pdf}
\vspace{-5mm}
\caption{\textbf{Yellow Legacy harness refinement is concentrated and recurrent rather than uniform.} (a)~Counts of CRUD operations (creation, update, delete) on skill and sub-agent definitions, binned per 2{,}000 turns. The harness is updated throughout the run rather than converging to a fixed scaffold. (b)~Update counts for the five most-updated components over the same horizon. A small subset of navigation and battle components accounts for the majority of updates.}
\label{fig:gpp_yellow}
\end{figure}

\begin{figure}[t]
\centering
\includegraphics[width=\linewidth]{h0_battle_complexity.pdf}
\vspace{-5mm}
\caption{Decision-making complexity of the Yellow Legacy \texttt{battle\_strategist\_agent} prompt at successive revisions during the Elite Four phase of the run: total nodes, decision gates, graph depth, and max fan-out. See appendix~\ref{app:gpp} for details.}
\label{fig:gpp_battle_complexity}
\end{figure}

\subsection{Continual Harness closes the gap to a hand-engineered harness}\label{sec:exp:closegap}

\begin{figure}[t]
\centering
\includegraphics[width=\linewidth]{h5_progression.pdf}
\vspace{-5mm}
\caption{Milestones reached vs.\ cumulative button presses. Red (left): 11-milestone subset sequence through Thunder Badge. Emerald (right): 9-milestone sequence through Knuckle Badge (2nd gym); x-axis capped at 8.5k. Lines stop at each run's last monitored milestone. Thick lines: seed medians; faint lines: individual seeds.}
\label{fig:progression}
\end{figure}

\cref{fig:progression} plots milestones reached against cumulative button presses for $\mathcal{H}_{\min}$, the three Continual Harness variants, and $\mathcal{H}_\mathrm{expert}$. On both games, Continual Harness substantially reduces the button-press cost of every monitored milestone relative to $\mathcal{H}_{\min}$ and recovers a majority of the $\mathcal{H}_{\min}$-to-expert efficiency gap, without access to the game decompilation, the milestone schedule, or any of the hand-built sub-agents that constitute $\mathcal{H}_\mathrm{expert}$. The residual gap to the expert harness concentrates in dialogue-heavy gym interiors and multi-turn battle strategy, components Continual Harness does not yet synthesize reliably; we attribute these to specific refinement targets in \cref{sec:exp:mechanisms}. On Red, the bootstrap-updating variant is more efficient than from-scratch at every milestone, indicating that the refinement signal compounds within the episode: a harness refined in a prior run accelerates the next even when the game state itself resets. Thus, automated refinement over harness components recovers a substantial fraction of the efficiency of a hand-engineered harness starting from a minimalist interface.
\subsection{Continual Harness gain depends on model capability}\label{sec:exp:substitute}

\begin{figure}[t]
\centering
\includegraphics[width=\linewidth]{h1_pareto.pdf}
\caption{Emerald cost--completion Pareto plane. Filled markers: individual 24-hour seeds. Ringed markers: per-cell medians. Dashed staircase: cost-monotone Pareto frontier. Y axis: fraction of the 31-milestone Emerald set reached; X axis: Gemini API spend (log scale, cached input at $25\%$).}
\label{fig:pareto}
\end{figure}

\cref{fig:pareto} compares every Emerald run from every model-harness cell with respect to cost and completion. On Pro, Continual Harness is strictly Pareto-dominant over $\mathcal{H}_{\min}$: from-scratch $\mathcal{H}_\mathrm{CH}$ reaches $100\%$ of milestones at a \$130 median, against $\mathcal{H}_{\min}$ at $98\%$ for \$215, a $\sim 40\%$ cost reduction with no completion loss. The two bootstrap variants on Pro reach $96$--$100\%$ of milestones at \$110--\$140. 
On Flash, harness benefit is high variance: bootstrap-updating reaches $80\%$ at \$42, marginally above $\mathcal{H}_{\min}$ at $77\%$ for \$30, while from-scratch and bootstrap-frozen variants have a higher variance. 
Flash-Lite with $\mathcal{H}_{\min}$ reaches $20\%$ at \$11; every Continual Harness variant on Flash-Lite falls to $3$--$13\%$ at comparable or higher cost. The harness gains requires a model that will sufficiently utilize the harness components properly.

\subsection{Open-source students co-learn with a refining harness}\label{sec:exp:colearn}

We test whether an open-source model improves its gameplay using the self-refining harness with reset-free training (batch size $= 1$). The model is first primed by supervised fine-tuning on frontier Continual Harness trajectories and an offline GRPO stage on a per-step process reward; neither warm-up stage produces meaningful milestone advancement on its own (Appendix~\ref{app:training}), and the live in-game gains we report here begin only at the co-learning stage. Each training iteration is a $K{=}256$-step DAgger~\citep{ross2011reduction,karten2026smallexperts} rollout through the full Continual Harness (memory, skills, sub-agents, and prompt all evolving via \cref{fig:methodology}), followed by a process-reward-model scoring pass, a Gemini-3.1-pro teacher relabel of low-reward windows, and a soft SFT update on the relabeled shard. The training loop is \emph{reset-free}: the emulator state at the end of iteration $k$ is loaded as the start of iteration $k{+}1$, so each curve in \cref{fig:colearn_pipeline} is a single agent's in-game trajectory traversed across its own training, not an aggregate over independent rollouts.

\begin{figure}[t]
\centering
\includegraphics[width=\linewidth]{analysis/figures/h6_colearn_pipeline.pdf}
\caption{\textbf{Reset-free DAgger+PRM training drives sustained milestone progress on Pok\'emon Red.} Milestone index reached versus training iteration $k$ for the five advancing runs; the broken y-axis labels each band's start and end. Filled dots: beginning of game. Open rings: mid-game checkpoint. Stars: judge-verified advances. $+N$: net objective gain. Dashed line: untrained Gemma-4 baseline (zero advance beyond the starting milestone). Teacher model: Gemini-3.1-pro.}
\label{fig:colearn_pipeline}
\end{figure}

\cref{fig:colearn_pipeline} shows that the model's live in-game position advances across training iterations on every plotted run, both from the beginning of the game and from mid-game checkpoints. Both staircase types share the same qualitative shape, indicating that the training signal that drives the model forward from the start of the game also drives it forward from advanced checkpoints; the training procedure is not specific to the early-game distribution. As a negative control, cross-family Qwen3.5 (27B, 35B) without the supervised warm-up stage produces parseable tool calls but cannot leave the starting area in a live rollout (Appendix~\ref{app:h6_full}), ruling out a rollout-protocol artifact. Together with the cross-checkpoint generalization, these results support the co-learning claim: an open-source model trained on data collected from its own play through a continually refining harness improves its in-game position iteration over iteration, without ever resetting the environment. Per-run identifiers, hyperparameters, and the per-iteration process-reward decomposition are reported in Appendix~\ref{app:training:dagger_resetfree}.

\subsection{Skills measurably self-improve toward an oracle}\label{sec:exp:mechanisms}

\begin{figure}[t]
\centering
\includegraphics[width=\linewidth]{h2_pathfinding.pdf}
\caption{Pathfinding skill mechanism. (Left)~Path-cost deficit of the top-10\% evolved navigation skill set (Gemini~3.1~Pro) relative to the Dijkstra oracle over a 24-hour run on warp-to-warp obstacle-navigation tasks; lower is better, dashed line at 0\% marks the oracle. (Right)~Cumulative navigation-skill invocations against button presses across the same conditions.}
\label{fig:pathfinding}
\end{figure}

We score refined navigation skills by their path cost relative to a Dijkstra oracle. This gives a direct measure of skill self-improvement, independent of end-task efficiency. \cref{fig:pathfinding} reports this measurement on warp-to-warp obstacle-aware navigation between fixed map entry and exit points, where greedy open-field hopping fails. Sub-agents, memory, prompt, and reset-free bootstrap transfer are deferred to Appendices~\ref{app:mechanisms} and~\ref{app:bootstrap}.

$\mathcal{H}_{\min}$ never invokes a navigation skill; every Continual Harness condition accumulates hundreds of invocations over a 24-hour run (\cref{fig:pathfinding}, right). On from-scratch runs the path-cost deficit falls from a near-half-cost penalty at the start to single digits early on and stays there (\cref{fig:pathfinding}, left). This improvement is in-loop and reset-free: failures from earlier invocations are diagnosed by the Refiner and the affected skills are repaired before later invocations within the same episode. Bootstrap-updating inherits a refined skill set and matches or outperforms bootstrap-frozen throughout, so continued refinement still adds value on top of an inherited set; bootstrap-frozen's flat trajectory bounds inheritance without further refinement.
